\section{A common truncation for all color constraints}
\label{sec:c-palettes}

The preceding two sections use input-list or omitted-set averages.
We now seek a single probability law satisfying bounds for every
coloring of either half of the labels. A small mean error for the next
card cannot be used unchanged after conditioning on a rare prescribed
mapping. Instead, a symmetry of the last $d-1$ layers and a random label
order will produce one event on switch settings where the accumulated
errors are small before any mapping constraint is selected.

We use positions $V=\mathbb F_2^d$, put $n=2^d$ and $h=n/2$, and let a
permutation send initial labels to final positions.  Layer $s$ uses
independent fair switches in direction $i_s=1+(s\bmod d)$.  A block of
$T=C_*d$ layers, with $C_*$ a positive integer, is a block of $C_*d$
physical shuffles because the rotating coordinate frame returns at its
endpoint.  Write $\mathbb P$ for its fair law on switch settings, and
$g$ for the endpoint permutation.  All constants below are fixed before
$d$ tends to infinity.

Fix an order $\omega=(c_1,\ldots,c_n)$ of the labels.  After revealing
the entire trajectories of $c_1,\ldots,c_r$, let $B_s$ be the available
positions at time $s$, namely those not occupied by these revealed labels,
and put $m=n-r$.  Let $p_s(x)$ be the conditional
probability that $c_{r+1}$ occupies $x$.  Define
\[
 w_s(x)=p_s(x)-m^{-1}{\bf1}_{B_s}(x),\qquad
 V_s=\sum_{x\in V}w_s(x)^2.
\]
Here $B_s$ and $w_s$ are the available set and centered vector denoted
$V_t$ and $f_t$ in Section~\ref{sec:first-route}; $V_s$ is their realized
squared energy, whose expectation is the quantity denoted $a_t$ there.
The conditional probabilities are still taken under the original fair law.
Lemma~\ref{lem:marginal-averaging} supplies the update: two available
input positions average their masses; one available input position
transports its mass to the unique available output; and an edge with no
available input carries zero.  Uniform mass on the available set obeys
the same rule.  Consequently $\sum_x w_s(x)=0$, $V_s\le1$, and $V_s$ is
nonincreasing.  The same lemma shows that, even after all revealed paths
are prescribed, the values at time $s$ depend only on their prefixes
through $s$.  This adaptedness will permit us to condition before a group
of layers; it does not redefine $p_s$ as a law conditioned on all past
switches.

\begin{lemma}[Independent translations in the two halves]
\label{lem:c-palette-half-translations}
Fix the label order and all switch settings through time
$a=s-d+1$, where $s\ge d-1$.  In each of the two halves determined by
coordinate $i_s$, independently translate the other $d-1$ coordinates.
The conditional distribution of $(B_s,w_s)$ is invariant under these two
translations.
\end{lemma}

\begin{proof}
The layers from $a$ through $s-1$ use each coordinate other than $i_s$
once and do not cross the two halves.  In either half build the desired
translation as those coordinates are processed.  If an accumulated
translation is already present, reindex the next layer's coins by that
translation; it preserves the matching.  If the current coordinate is
to be toggled, complement every coin of this layer in that half.  This
is a bijection of the fair coin arrays and translates all output
positions in that coordinate.  At an edge with one available input the
mass moves to the translated available output.  At an edge with two
available inputs both outputs receive the same average, so complementing the coin has the
same equivariance.  Induction over the $d-1$ layers proves the assertion,
and the transformations in the two halves use disjoint coin arrays.
\end{proof}

The symmetry makes a matching behave, in expectation, like the complete
bipartite graph between the halves.  This is useful even though the
available positions and their masses in a given half are dependent.

\begin{proposition}[Mean endpoint discrepancy]
\label{prop:c-palette-discrepancy}
For every fixed $0<\delta<1/4$, there is a positive integer $C_*$ such
that the following holds for all sufficiently large $n=2^d$.  Choose
$\omega$ uniformly and independently of the switches, let $T=C_*d$, and
put
\[
 \Delta_r=\max_{x\in B_T}|p_T(x)-m^{-1}|.
\]
Then, uniformly for $n-r=m\ge\delta n$,
\begin{equation}
 \mathbb E_{\omega,\mathbb P}\Delta_r\le n^{-30}.
 \label{eq:c-palette-discrepancy}
\end{equation}
\end{proposition}

\begin{proof}
Immediately before layer $s$, put $H_b=\{x\in V:x_{i_s}=b\}$ and
$m_b=|B_s\cap H_b|$ for $b=0,1$. Call the configuration balanced if
$m_0,m_1\ge m/4$.  The energy lost at an edge $\{x,y\}$ with two
available positions is $(w_s(x)-w_s(y))^2/2$.  The balance event is translation invariant.
Averaging the translation of one half in
Lemma~\ref{lem:c-palette-half-translations}, the sum over matching edges
becomes $1/h$ times the sum over all pairs of available positions in
opposite halves.  The latter sum equals
\[
 m_0\sum_{x\in H_1}w_s(x)^2+
 m_1\sum_{x\in H_0}w_s(x)^2
 -2\left(\sum_{x\in H_0}w_s(x)\right)
    \left(\sum_{x\in H_1}w_s(x)\right).
\]
The last term is nonnegative because the two sums add to zero.  On
balance the expression is therefore at least $(m/4)V_s$.

Conditional on the complete switch setting, the last $m$ labels of the
random order form a uniform $m$-set, whose image at time $s$ is again a
uniform $m$-set.  Each half-count has mean $m/2$.  Exposing the set by
sampling without replacement gives a Doob martingale with increments
of absolute value at most one: couple two possible next draws by
exchanging the drawn elements in the remaining completion.  The
bounded-increment exponential-moment estimate gives
\[
 \mathbb P\{m_b<m/4\}\le \exp(-m/32),\qquad b=0,1.
\]
This is an unconditional bound for the joint experiment.  No
concentration conditional on the revealed trajectories is required.
Since $V_s\le1$, for $s\ge d-1$ we obtain
\[
 \mathbb EV_{s+1}
 \le\left(1-\frac{m}{8h}\right)\mathbb EV_s
       +O(e^{-c_\delta n})
 \le(1-\delta/4)\mathbb EV_s+O(e^{-c_\delta n}).
\]
The symmetry windows may overlap: each use gives an expectation
inequality and asserts no independence between windows.  Iteration,
followed by $\Delta_r\le\sqrt{V_T}$ and Jensen's inequality, gives
\[
 \mathbb E\Delta_r
 \le\left((1-\delta/4)^{T-d+1}
                  +O(e^{-c_\delta n})\right)^{1/2}.
\]
Choosing the integer $C_*$ sufficiently large proves
\eqref{eq:c-palette-discrepancy}, uniformly in $m$.
\end{proof}

\subsection{Selecting the partition before selecting a color constraint}

Put $v=\lfloor\delta n\rfloor$.  First choose a uniform partition
$V=A\sqcup B$ with $|A|=|B|=h$.  There are two order distributions:
$A$ first and then $B$, each uniformly ordered internally, and the
reversed convention.  Before the partition is fixed, each distribution
is a uniform order of all labels.  Proposition~\ref{prop:c-palette-discrepancy}
therefore allows us to choose a deterministic partition for which
\[
 \mathbb E_{\mathbb P}Z\le2n^{-29},\qquad
 Z=\mathbb E_{\omega:A,B}\sum_{r=0}^{n-v-1}\Delta_r
   +\mathbb E_{\omega:B,A}\sum_{r=0}^{n-v-1}\Delta_r.
\]
The last retained rank leaves at least $v+1>\delta n$ sites, as needed.
Fix this partition from now on, and retain the event
\[
 \mathcal E=\{Z\le n^{-6}\},\qquad z=\mathbb P(\mathcal E).
\]
Markov's inequality gives
\begin{equation}
 1-z\le2n^{-23}.
 \label{eq:c-palette-retained-mass}
\end{equation}
Let $\mu$ be the endpoint law under $\mathbb P$ and $\nu$ its endpoint
law under $\mathbb P(\,\cdot\mid\mathcal E)$.  The event $\mathcal E$
lives on switch settings and need not be determined by the endpoint.
Nevertheless data processing gives
$\|\nu-\mu\|_{\rm TV}\le1-z$.

A palette of $B$ is a partition into nonempty color classes of sizes
$l_1,\ldots,l_u$.  Set $p_i=l_i/h$ and
$\mathcal H(p)=-\sum_i p_i\log p_i$, and let $H$ permute each color
class while fixing $A$ pointwise.  The coset $gH$ fixes every output of
$A$ and the output set of each color class of $B$.

We apply the information lemma for retained reveal experiments
\cite[Lemma~\ref*{l-l:palette-information}]{partialFourier},
\emph{with both order distributions and this same event $\mathcal E$}.
Its hypothesis is a pointwise cap on the
order-average sum of endpoint errors, together with the fair conditional
bound
\[
 \mathbb P\{\hbox{next endpoint lies in a target set of size }a
               \mid\hbox{revealed paths}\}
 \le a/m+n\Delta_r.
\]
Here the latter follows by summing the conditional probabilities over
the available target sites, and the former is exactly the definition
of $\mathcal E$. The role of the pointwise bound is seen by conditioning
the fair law on $\mathcal E\cap\{g\in g_0H\}$, for a fixed coset
$g_0H$ when this event has positive probability. The lemma's entropy chain gives
\[
 -\log\mathbb P(\mathcal E\cap\{g\in g_0H\})
 \ge \log[S_n:H]-v\mathcal H(p)-n^2n^{-6}.
\]
At each revealed rank, a prescribed target costs its uniform
logarithmic probability up to $n^2\Delta_r$. Their order-averaged sum
is bounded on every retained setting, so the same error bound holds
under this rare conditioning. The missing $v$ ranks belong to the
second half; their mean logarithmic multinomial count is at most
$v\mathcal H(p)$. Dividing the resulting probability bound by $z$ gives
\begin{equation}
 \nu(gH)\le [S_n:H]^{-1}
     \exp\{v\mathcal H(p)+n^{-4}-\log z\},
 \label{eq:c-palette-coset-cap}
\end{equation}
simultaneously for all palettes and all cosets, and also with $A,B$
interchanged. The mean-discrepancy estimate
\eqref{eq:c-palette-discrepancy} was used to choose $\mathcal E$ under
the original law; only the pointwise bound on $\mathcal E$ enters
the rare-event argument.

\begin{corollary}[Palette application]
\label{cor:c-palette-mixing}
There is an absolute positive integer $C_*$ for which the full Thorp
permutation law after $30C_*d$ physical shuffles has total-variation
distance tending to zero from uniform, uniformly over starting decks.
\end{corollary}

\begin{proof}
The signed-palette lemma
\cite[Lemma~\ref*{l-l:signed-palette}]{partialFourier} gives an absolute
constant $C_{\rm pal}$: every irreducible half-group type $\beta$ has a
palette and a product of trivial or sign characters occurring in it,
with $h\mathcal H(p)\le C_{\rm pal}\log D_\beta$.
Choose $\delta<1/4$ so that $2\delta C_{\rm pal}\le1/10$, and only
then choose $C_*$ in Proposition~\ref{prop:c-palette-discrepancy}.
Since $v\le2\delta h$ and $n^{-4}-\log z=o(1)$,
\eqref{eq:c-palette-coset-cap} gives
\[
 \nu(gH)\le \frac{2D_\beta^{1/10}}{[S_n:H]}
\]
for large $n$.  It holds for every placement of these color classes
and hence every conjugate in either half-group.  More precisely,
\eqref{eq:c-palette-coset-cap}, with
$\varepsilon_n=n^{-4}-\log z=o(1)$ and the chosen $\delta$, verifies
the full entropy-form hypotheses for every palette in the signed-palette
transfer \cite[Theorem~\ref*{l-l:palette-transfer}]{partialFourier}, whose
conclusion is
\[
 \|\widehat\nu(\lambda)\|_{\rm op}\le D_\lambda^{-1/6}
 \quad\bigl(\lambda\ne(n),(1^n)\bigr).
\]
Here $D_\lambda$ is the irreducible degree and
$\widehat\nu(\lambda)=\sum_g\nu(g)\rho_\lambda(g)$ in a unitary model.

The original block has zero mean sign: conditioning on all but one
coin in its final layer, flipping that coin composes the endpoint with
a transposition.  Clipping on settings therefore gives
$|\widehat\nu(\mathrm{sgn})|\le(1-z)/z=o(1)$.
For thirty independent factors, Plancherel and
$\|M\|_{\rm HS}\le\sqrt{D_\lambda}\|M\|_{\rm op}$ bound the
nontrivial, nonsign contribution by
\[
 \sum_{\lambda\ne(n),(1^n)}D_\lambda^{2-30/3}
 =\sum_{\lambda\ne(n),(1^n)}D_\lambda^{-8}=o(1),
\]
using the reciprocal-degree conclusion proved in that theorem.  The sign contribution
also tends to zero.  Thus $\nu^{*30}$ tends to uniform in total
variation.  Replacing its thirty factors by $\mu$ costs at most
$30(1-z)=o(1)$.  Each factor has $C_*d$ physical shuffles, and changing
the starting deck translates the group law.  This proves the claim.
\end{proof}
